\documentclass[11pt,a4paper]{article}
\usepackage[T1]{fontenc}
\usepackage[utf8]{inputenc}
\usepackage{newtxtext}
\usepackage{amsmath,amsthm,mathtools}
\usepackage{newtxmath}
\usepackage[margin=27mm,headheight=14pt]{geometry}
\usepackage{microtype,booktabs,tabularx,array}
\usepackage[dvipsnames]{xcolor}
\usepackage{enumitem,fancyhdr,etoolbox,needspace}
\usepackage{hyperref}
\usepackage{bookmark}
\hypersetup{colorlinks=true,linkcolor=MidnightBlue,citecolor=MidnightBlue,
 urlcolor=MidnightBlue,pdftitle={Periods, Global Fueter Primitives, and Integrable Spherical Singularities},
 pdfauthor={OpenAI},pdfsubject={Quaternionic analysis; global inverse Fueter problem; residues and energy}}
\urlstyle{same}
\numberwithin{equation}{section}
\newtheorem{theorem}{Theorem}[section]
\newtheorem{proposition}[theorem]{Proposition}
\newtheorem{lemma}[theorem]{Lemma}
\newtheorem{corollary}[theorem]{Corollary}
\theoremstyle{definition}
\newtheorem{definition}[theorem]{Definition}
\newtheorem{example}[theorem]{Example}
\theoremstyle{remark}
\newtheorem{remark}[theorem]{Remark}
\newcommand{\R}{\mathbb R}
\newcommand{\C}{\mathbb C}
\newcommand{\HH}{\mathbb H}
\newcommand{\HC}{\mathbb H_{\C}}
\newcommand{\SSS}{\mathbb S}
\newcommand{\ii}{\mathbf i}
\newcommand{\jj}{\mathbf j}
\newcommand{\kk}{\mathbf k}
\newcommand{\ic}{\iota}
\newcommand{\D}{\mathcal D}
\newcommand{\lap}{\Delta_4}
\newcommand{\SR}{\mathcal{SR}}
\newcommand{\AM}{\mathcal{AM}}
\newcommand{\Aff}{\mathcal{A}_{\mathrm{aff}}}
\newcommand{\Per}{\operatorname{Per}}
\newcommand{\Ind}{\operatorname{Ind}}
\newcommand{\Rea}{\operatorname{Re}}
\newcommand{\Ima}{\operatorname{Im}}
\newcommand{\supp}{\operatorname{supp}}
\newcommand{\dd}{\,\mathrm d}
\newcommand{\norm}[1]{\lVert #1\rVert}
\newcommand{\abs}[1]{\lvert #1\rvert}
\newcommand{\cI}{\mathcal I}
\newcommand{\Gz}{G_{p,0}}
\newcommand{\Go}{G_{p,1}}
\newcommand{\Sp}{S_p}
\setlength{\emergencystretch}{2em}
\setlength{\parskip}{2pt}
\setlist{itemsep=2pt,topsep=4pt}
\pagestyle{fancy}
\fancyhf{}
\fancyhead[L]{\small Periods and spherical singularities in quaternionic analysis}
\fancyhead[R]{\small\thepage}
\renewcommand{\headrulewidth}{0.3pt}
\title{\vspace{-12mm}\textbf{Periods, Global Fueter Primitives,\\and Integrable Spherical Singularities}\\[6pt]
\large A constructive study in quaternionic axial analysis}
\author{Research manuscript prepared with OpenAI}
\date{21 September 2026}
\begin{document}
\maketitle
\thispagestyle{empty}
\begin{abstract}
The local inverse Fueter theorem does not, by itself, decide whether a
single-valued slice-regular primitive exists on a multiply connected domain.
We construct two closed quaternion-valued one-forms directly from an axial
Fueter-regular function and prove that their periods are the complete global
obstruction. For a meridian with $m$ holes, the obstruction map is onto
$\HH^{2m}$ and has an explicit splitting; this yields a period-corrected Runge
approximation theorem and quantitative obstructions to approximation without
correction. On domains meeting the real axis, only the conjugation-invariant
part of first cohomology contributes: a conjugate pair of holes contributes
two quaternionic obstructions, whereas a conjugation-fixed hole contributes
none. We then classify every locally $L^1$ axial Fueter-regular singularity
along a nonreal conjugacy sphere by two quaternionic residues, modulo a
removable function. These residues determine both an affine-density
Cauchy--Fueter source supported on the sphere and the exact coefficient of
logarithmic $L^2$-energy divergence. A pair of conserved three-dimensional
boundary fluxes recovers both residues. In particular, $L^2$ implies removability;
below $L^1$, vanishing residues no longer suffice. All new assertions are
proved, and deterministic symbolic and numerical checks are supplied for the
explicit formulas. The local inverse theorem and the underlying spherical
kernels are established results, not claimed as new here.
\end{abstract}
\medskip
\noindent\textbf{Keywords.} Inverse Fueter problem; slice regularity;
axial monogenic functions; monodromy; quaternionic residues; removability;
renormalized energy; Runge approximation.
\medskip
\noindent\textbf{Mathematical scope.} One quaternionic variable; left regularity;
coefficients on the right. The global cokernel calculation is constructive for
finite planar topology. The period criterion itself requires no finite-topology
hypothesis. The singularity theorems concern complete nonreal conjugacy spheres,
not arbitrary singular sets or arbitrary nonaxial functions.
\newpage
\begingroup\small
\setlength{\parskip}{0pt}
\makeatletter\patchcmd{\l@section}{1.0em}{0.45em}{}{}\makeatother
\tableofcontents
\endgroup
\newpage

\section{The questions and the contribution}
\label{sec:intro}

The Fueter map sends a slice-regular quaternionic function $f$ to the
Fueter-regular function $\lap f$. Its kernel consists of affine functions
$q\mapsto qa+b$, and axial Fueter-regular functions have local slice-regular
primitives. These facts belong to established inverse Fueter theory
\cite{CSS,CPSS}. The supplied expository manuscript \cite{Supplied} ends its
inverse construction with a simply connected meridian and explicitly notes
that periods may prevent global inversion. We take that observation as a
starting point, not as a theorem to be rediscovered.

There are three concrete questions. First, can the global obstruction be read
from the target function itself, without choosing successive local primitives?
Second, which obstruction classes actually occur, and how do they affect
approximation by Fueter images? Third, near an isolated singular sphere, can
the same obstruction data classify the singularity rather than merely detect
failure of a primitive?

The answers developed below are linked by a pair of quaternionic numbers per
cycle. For $g(x+Ir)=P(x,r)+I Q(x,r)$, define
\begin{align}
 \omega_g&=\frac12(-P\dd x+Q\dd r),\label{eq:intro-omega}\\
 \eta_g&=\frac12\bigl((xP+rQ)\dd x+(rP-xQ)\dd r\bigr).
 \label{eq:intro-eta}
\end{align}
Both forms are closed. The important feature of $\eta_g$ is that it contains
\emph{only} $P$ and $Q$: no previously chosen primitive appears in it.
For a closed meridian curve $\gamma$, put
\[
 a_\gamma=\int_\gamma\omega_g,\qquad
 b_\gamma=\int_\gamma\eta_g.
\]
A local primitive continued around $\gamma$ changes by $qa_\gamma+b_\gamma$.
Theorem~\ref{thm:period} proves that there is no additional obstruction.

For a puncture $p=u+\ic v$, $v>0$, its circularization is the sphere
\[
 S_p=\{u+Iv:I\in\SSS\}.
\]
Our strongest local conclusion is that, under an $L^1$ hypothesis, the two
periods $a,b$ around $p$ determine the entire nonremovable part. Its
four-dimensional Dirac defect is
\begin{equation}
 \D g=\frac{2}{v}(qa+b)\,\delta_{S_p},
 \label{eq:intro-defect}
\end{equation}
where $\delta_{S_p}$ is surface measure, not normalized probability measure.
If $\rho(q)=\bigl((\Rea q-u)^2+(|\Ima q|-v)^2\bigr)^{1/2}$, then
\begin{equation}
 \int_{\varepsilon<\rho(q)<R}|g(q)|^2\dd V(q)
 =8\bigl(|ua+b|^2+v^2|a|^2\bigr)\log\frac R\varepsilon+O(1).
 \label{eq:intro-energy}
\end{equation}
In fact, subtracting the displayed logarithm leaves a finite limit.
Thus monodromy, a surface-supported source, and energy concentration encode
exactly the same eight real parameters.

\subsection{What is, and is not, claimed as new}
The precise results proposed as contributions are the direct two-form period
criterion and its explicit finite-hole splitting; its real-axis symmetry
version; the period-corrected approximation statement; and, most importantly,
the complete $L^1$ normal form together with its exact defect and energy
identities. Proofs, rather than numerical evidence, establish these statements
within this manuscript.

The local inverse, the affine kernel, holomorphic stem functions, and
spherical kernels obtained from arctangents are not new. In particular,
Colombo--Pe\~na Pe\~na--Sabadini--Sommen already compute spherical kernels and
arctangent primitives \cite[Example 2]{CPSS}. We use equivalent local
logarithmic representatives and give their period normalization. Recent work
on factorized Fueter maps and polyanalytic extensions \cite{DP2026} is also
part of the relevant background.

A targeted literature check, through 21 September 2026, did not locate the
exact combined statements proved here. This is a qualified assessment of
novelty, not a certification of priority or a claim to have resolved a named
published conjecture. In particular, general sheaf methods make the existence
of some cohomological obstruction natural. The more specific content here is
the explicit pair of forms, complete realization of the classes, the
integrability threshold, and the normalized relation between periods, defects,
and energy. Section~\ref{sec:literature} records the comparison and its limits.

\subsection{Prerequisites and conventions}
We use elementary complex analysis, including its identity theorem, Cauchy
estimates, and Runge approximation; elementary homology of finitely connected
planar domains; and the divergence theorem. The weak removable-singularity
lemma needed at the $L^1$ endpoint is proved below. Every quaternionic
assertion on which the new results depend is either proved directly or reduced
to these scalar facts. No assertion is made about unrestricted quaternionic
pointwise products preserving regularity.

\section{Axial functions and the local differential identities}
\label{sec:prelim}

\subsection{Quaternionic and complexified notation}
Write
\[
 q=x_0+\ii x_1+\jj x_2+\kk x_3,
 \qquad \ii^2=\jj^2=\kk^2=\ii\jj\kk=-1.
\]
Quaternionic conjugation is denoted by $\bar q$ and $|q|^2=q\bar q$.
Let
\[
 \SSS=\{I\in\HH:I^2=-1\},\qquad
 \D=\partial_{x_0}+\ii\partial_{x_1}+\jj\partial_{x_2}+\kk\partial_{x_3}.
\]
The conjugate operator satisfies $\overline\D\D=\D\overline\D=\lap$.
The central complex unit $\ic$ in
\[
 \HC=\HH\otimes_\R\C=\{A+\ic B:A,B\in\HH\}
\]
is not a quaternionic imaginary unit. Its conjugation is
$\kappa(A+\ic B)=A-\ic B$. In particular, $\kappa$ does not conjugate $A$ or $B$.
Use the Euclidean norm $|A+\ic B|_c^2=|A|^2+|B|^2$.

Let $U$ be a connected open subset of $\C^+=\{x+\ic r:r>0\}$. Its
circularization is
\[
 \Omega_U=\{x+Ir:x+\ic r\in U,\ I\in\SSS\}.
\]
A holomorphic function $F=A+\ic B:U\to\HC$ induces
\[
 f=\cI(F),\qquad f(x+Ir)=A(x,r)+I B(x,r).
\]
Equivalently, reflect $F$ to $U\cup\bar U$ by
$F(\bar z)=\kappa(F(z))$ and use the usual stem construction
\cite{GP}. We call these functions slice regular and write $\SR(U)$.
This definition on a product domain is stronger than an unrelated holomorphic
choice on every quaternionic slice.

An axial function $g=P+I Q$ has coefficients $P,Q:U\to\HH$. We write
$\AM(U)$ for the smooth axial functions annihilated by $\D$. All these spaces
are right $\HH$-modules. They are not being regarded as left $\HH$-modules.

\subsection{Axial calculus}
\begin{lemma}[Axial formulas]\label{lem:axial}
For a smooth axial function $h(x+Ir)=A(x,r)+I B(x,r)$,
\begin{align}
 \D h&=A_x-B_r-\frac{2B}{r}+I(B_x+A_r),\label{eq:axial-D}\\
 \lap h&=A_{xx}+A_{rr}+\frac{2A_r}{r}
 +I\left(B_{xx}+B_{rr}+\frac{2B_r}{r}-\frac{2B}{r^2}\right).
 \label{eq:axial-lap}
\end{align}
Consequently, $g=P+I Q$ belongs to $\AM(U)$ exactly when
\begin{equation}
 P_x-Q_r-\frac{2Q}{r}=0,\qquad Q_x+P_r=0.
 \label{eq:vekua}
\end{equation}
For $f=\cI(A+\ic B)\in\SR(U)$,
\begin{equation}
 A_x=B_r,\quad A_r=-B_x,\qquad
 \lap f=\frac{2A_r}{r}
  +I\left(\frac{2B_r}{r}-\frac{2B}{r^2}\right).
 \label{eq:fueter-formula}
\end{equation}
\end{lemma}
\begin{proof}
Put $\boldsymbol v=\Ima q$ and $I=\boldsymbol v/r$. For $\ell=1,2,3$,
\[
 \partial_{x_\ell}r=\frac{x_\ell}{r},\qquad
 \partial_{x_\ell}I=\frac{e_\ell}{r}-\frac{\boldsymbol v x_\ell}{r^3},
 \quad(e_1,e_2,e_3)=(\ii,\jj,\kk).
\]
It follows that $\sum e_\ell\partial_{x_\ell}I=-2/r$,
$\Delta_{\R^3}I=-2I/r^2$, and
$\sum(\partial_{x_\ell}I)(\partial_{x_\ell}r)=0$.
The product rule gives \eqref{eq:axial-D} and \eqref{eq:axial-lap}.
Evaluating an expression $C+I H$ at $I$ and $-I$ proves that it vanishes for
all $I$ exactly when $C=H=0$, which gives \eqref{eq:vekua}. The
Cauchy--Riemann equations for $A+\ic B$ imply that $A$ and $B$ are planar
harmonic and yield \eqref{eq:fueter-formula}.
\end{proof}

\begin{proposition}[Fueter map]\label{prop:fueter}
The map $T:\SR(U)\to\AM(U)$ given by $Tf=\lap f$ is well defined.
\end{proposition}
\begin{proof}
For the two coefficients in \eqref{eq:fueter-formula}, differentiation gives
\[
 \partial_x\left(\frac{2A_r}{r}\right)
 =\partial_r\left(\frac{2B_r}{r}-\frac{2B}{r^2}\right)
   +\frac2r\left(\frac{2B_r}{r}-\frac{2B}{r^2}\right),
\]
using $A_x=B_r$. The other equation in \eqref{eq:vekua} follows from
$A_r=-B_x$ and its derivatives. Thus $\D(Tf)=0$.
\end{proof}

\subsection{Two potentials instead of successive primitives}
\begin{lemma}[Potential identities]\label{lem:potentials}
Let $Tf=g=P+I Q$, where $f=\cI(A+\ic B)$. Set
\begin{equation}
 C=\frac Br,\qquad D=A-x\frac Br.
 \label{eq:potentials}
\end{equation}
Then
\begin{equation}
 \dd C=\omega_g,\qquad \dd D=\eta_g,
 \qquad f(x+Ir)=(x+Ir)C(x,r)+D(x,r).
 \label{eq:potential-identities}
\end{equation}
Conversely, if $C,D$ satisfy the first two identities, the last expression
is slice regular and its Laplacian is $g$.
\end{lemma}
\begin{proof}
By \eqref{eq:fueter-formula} and the Cauchy--Riemann equations,
\[
 C_x=-\frac P2,\qquad C_r=\frac Q2.
\]
Furthermore, $A_x=B_r=C+rQ/2$ and $A_r=rP/2$. Hence
\[
 D_x=A_x-C-xC_x=\frac{xP+rQ}{2},\qquad
 D_r=A_r-xC_r=\frac{rP-xQ}{2}.
\]
For the converse, put $A=D+xC$ and $B=rC$. The displayed equations imply
\[
 A_x=C+\frac{rQ}{2}=B_r,\qquad A_r=\frac{rP}{2}=-B_x.
\]
Thus $A+\ic B$ is holomorphic, and substitution in
\eqref{eq:fueter-formula} gives $Tf=P+I Q$.
\end{proof}

The change from $A$ to $D=A-xC$ is the essential simplification. In a
successive-integration construction, the second one-form depends on the first
primitive. Here both obstruction forms are globally defined before any
integration is performed.

\section{The complete global period obstruction}
\label{sec:periods}

\begin{lemma}[Closedness]\label{lem:closed}
For every $g\in\AM(U)$, the forms $\omega_g$ and $\eta_g$ in
\eqref{eq:intro-omega}--\eqref{eq:intro-eta} are closed.
\end{lemma}
\begin{proof}
For $\omega_g$, equality of the cross derivatives is
$Q_x=-P_r$, the second equation of \eqref{eq:vekua}. For $\eta_g$, compute
\begin{align*}
 \partial_r(xP+rQ)&=xP_r+Q+rQ_r,\\
 \partial_x(rP-xQ)&=rP_x-Q-xQ_x
 =rQ_r+Q+xP_r.
\end{align*}
The last equality uses both equations in \eqref{eq:vekua}. All coefficients
are real scalar multiples of quaternions, so no multiplication order is
changed.
\end{proof}

Integration is componentwise in the real basis $1,\ii,\jj,\kk$. Closedness
and Stokes' theorem make
\begin{equation}
 \Per_g:H_1(U;\mathbb Z)\longrightarrow\HH\oplus\HH,
 \qquad [\gamma]\longmapsto
 \left(\int_\gamma\omega_g,\int_\gamma\eta_g\right)
 \label{eq:period-map}
\end{equation}
a well-defined homomorphism. Equivalently, the obstruction is the pair
$([\omega_g],[\eta_g])$ of real de Rham cohomology classes with quaternionic
coefficients.

\begin{theorem}[Global inverse and monodromy]\label{thm:period}
Let $U\subset\C^+$ be any connected open set and $g\in\AM(U)$.
The following are equivalent:
\begin{enumerate}[label=\textup{(\roman*)}]
\item $g=\lap f$ for a single-valued $f\in\SR(U)$;
\item $\Per_g=0$;
\item both $\omega_g$ and $\eta_g$ are exact on $U$.
\end{enumerate}
When these conditions hold, fix $z_0\in U$ and $c,d\in\HH$. Every primitive
is obtained from
\begin{equation}
 f(x+Ir)=(x+Ir)\left(c+\int_{z_0}^{x+\ic r}\omega_g\right)
       +d+\int_{z_0}^{x+\ic r}\eta_g.
 \label{eq:inverse-formula}
\end{equation}
The integrals are independent of the path. The ambiguity is precisely
$q\mapsto qa+b$, $a,b\in\HH$.

Without assuming vanishing periods, the same formula on path classes gives a
primitive on the universal cover. Continuation around $\gamma$ changes it by
\begin{equation}
 f_\gamma(q)-f(q)=q a_\gamma+b_\gamma.
 \label{eq:monodromy}
\end{equation}
\end{theorem}
\begin{proof}
If $g=Tf$, Lemma~\ref{lem:potentials} gives global potentials, so (i) implies
(iii). Exact forms have zero integrals on closed paths. Conversely, if a
closed form has zero periods, integration from $z_0$ along any piecewise
smooth path is path independent and produces a primitive; this argument
applies to all eight real components. Thus (ii) implies (iii).
Lemma~\ref{lem:potentials} now proves (i) and \eqref{eq:inverse-formula}.

For $Tf=0$, the same lemma says that $C$ and $D$ have zero differential.
They are constant because $U$ is connected, giving $f(q)=qa+b$.
On the universal cover the pulled-back forms are exact. The change of their
potentials under a deck transformation is their integral around the
corresponding loop. Substitution in $f=qC+D$ proves \eqref{eq:monodromy}.
\end{proof}

\begin{remark}[The coefficient field matters]
The affine ambiguity is $\HH^2$, not $\HC^2$. For example, a stem that is
constant and equal to the central unit $\ic$ on $U$ induces $f(x+Ir)=I$,
whose Laplacian is $-2I/r^2$, not zero. Losing this distinction would double
the claimed cokernel incorrectly.
\end{remark}

\begin{corollary}[Abelian monodromy]\label{cor:abelian}
Every commutator loop acts trivially on a local Fueter primitive. The
obstruction depends only on first homology, not on the full nonabelian
fundamental group.
\end{corollary}
\begin{proof}
Both period integrals are homomorphisms into the additive abelian group
$\HH^2$. Their values on commutators vanish, and
\eqref{eq:monodromy} applies.
\end{proof}

\section{Realization of every finite-hole obstruction}
\label{sec:splitting}

\subsection{Explicit elementary generators}
For $p=u+\ic v\in\C^+$, let
\[
 \ell_p(z)=\frac{1}{2\pi\ic}\log(z-p)
\]
on any local branch. Define locally
\begin{equation}
 G_{p,0}=T\cI(\ell_p),\qquad
 G_{p,1}=T\cI(z\ell_p).
 \label{eq:generators}
\end{equation}
On changing a logarithm branch by $2\pi\ic n$, the two primitives change
by $n$ and $nq$. Their Laplacians therefore do not change. Thus the $G_{p,k}$
are globally defined elements of $\AM(\C^+\setminus\{p\})$.

For computation, put
\[
 X=x-u,\quad Y=r-v,\quad s=X^2+Y^2,\quad L=\frac12\log s.
\]
Writing $G_{p,k}=P_k+I Q_k$ gives the branch-free formulas
\begin{align}
 P_0&=\frac{X}{\pi r s},&
 Q_0&=-\frac{Y}{\pi r s}+\frac{L}{\pi r^2},\label{eq:G0}\\
 P_1&=\frac1{\pi r}\left(L+\frac{xX+rY}{s}\right),&
 Q_1&=\frac{xv-ru}{\pi r s}+\frac{xL}{\pi r^2}.
 \label{eq:G1}
\end{align}
Indeed, if $\theta$ is a local argument of $z-p$, then
\[
 \ell_p=\frac{\theta}{2\pi}-\ic\frac L{2\pi},\qquad
 z\ell_p=\frac{x\theta+rL}{2\pi}
          +\ic\frac{r\theta-xL}{2\pi}.
\]
Using $\theta_r=X/s$ and $L_r=Y/s$ in
\eqref{eq:fueter-formula} proves \eqref{eq:G0}--\eqref{eq:G1}.

\begin{proposition}[Normalization]\label{prop:normalization}
For every closed curve $\gamma\subset\C^+\setminus\{p\}$,
\begin{align}
 \Per_{G_{p,0}}([\gamma])&=(0,\Ind(\gamma,p)),\\
 \Per_{G_{p,1}}([\gamma])&=(\Ind(\gamma,p),0).
 \label{eq:normalized-periods}
\end{align}
\end{proposition}
\begin{proof}
The changes of the local primitives are respectively
$\Ind(\gamma,p)$ and $q\Ind(\gamma,p)$.
Theorem~\ref{thm:period} identifies their coefficients with the periods.
Uniqueness of these coefficients follows by evaluating $qa+b$ on an open set.
\end{proof}

\subsection{The split exact sequence}
For an explicit finite-topology statement, assume that $U$ is bounded,
$\overline U\subset\C^+$, and its complement in the Riemann sphere has
components $K_0,K_1,\ldots,K_m$, with $\infty\in K_0$.
Assume the usual finite planar topology: positively oriented curves
$\gamma_1,\ldots,\gamma_m\subset U$ form a homology basis, and points
$p_j\in K_j$ can be chosen so that
\begin{equation}
 \Ind(\gamma_i,p_j)=\delta_{ij}.
 \label{eq:winding-basis}
\end{equation}
For example, a bounded domain with finitely many disjoint smooth boundary
components, with finitely many additional punctures, satisfies these
hypotheses. Here ``positive'' means counterclockwise around a hole, not the
induced orientation of the boundary of $U$.

Let $\Aff(U)=\{q\mapsto qa+b:a,b\in\HH\}$. Package the periods as
\[
 \Pi_U(g)=\bigl((a_1,b_1),\ldots,(a_m,b_m)\bigr)\in\HH^{2m}.
\]

\begin{theorem}[Constructive cokernel calculation]\label{thm:split}
Under the preceding hypotheses there is an exact sequence of right
$\HH$-modules
\begin{equation}
 0\longrightarrow\Aff(U)\longrightarrow\SR(U)
 \xrightarrow{\;T\;}\AM(U)
 \xrightarrow{\;\Pi_U\;}\HH^{2m}\longrightarrow0.
 \label{eq:exact}
\end{equation}
The last map has the explicit right inverse
\begin{equation}
 S\bigl((a_j,b_j)_{j=1}^m\bigr)
 =\sum_{j=1}^m\bigl(G_{p_j,1}a_j+G_{p_j,0}b_j\bigr).
 \label{eq:split-map}
\end{equation}
In particular,
\begin{equation}
 \AM(U)=T\SR(U)\oplus
 \operatorname{span}_{\HH}\{G_{p_j,1},G_{p_j,0}:1\le j\le m\},
 \label{eq:direct-sum}
\end{equation}
and the cokernel has quaternionic dimension $2m$, or real dimension $8m$.
\end{theorem}
\begin{proof}
Exactness through $\AM(U)$ follows from Theorem~\ref{thm:period}, because the
chosen curves form a homology basis. Equation~\eqref{eq:normalized-periods}
and \eqref{eq:winding-basis} show that $\Pi_U S$ is the identity, proving
surjectivity and the splitting. For any $g$, the function
$g-S\Pi_Ug$ has zero periods and hence is a Fueter image. If a linear
combination of the displayed generators is a Fueter image, applying $\Pi_U$
forces every coefficient to be zero. This proves the direct sum and the
dimension statement.
\end{proof}

\begin{remark}[What is canonical]
The cohomology-valued obstruction is canonical. A coordinate vector for it
requires cycles; its splitting requires representatives $p_j$ and a choice of
generator within its Fueter-image coset. The direct-sum complement is
therefore not canonical. No finite-dimensionality assertion is made here for
infinitely connected meridians, although Theorem~\ref{thm:period} still applies.
\end{remark}

\begin{example}[Two independent obstructions on an annulus]
Take $U=\{z: r_0<|z-p|<r_1\}$ with $0<r_0<r_1<v$.
Then $G_{p,0}$ has $(a,b)=(0,1)$ and $G_{p,1}$ has $(a,b)=(1,0)$.
Neither has a single-valued slice-regular primitive on $\Omega_U$.
Moreover, testing only $\omega_g$ would incorrectly declare $G_{p,0}$
unobstructed. Both forms are necessary.
\end{example}

\section{Domains meeting the real axis}
\label{sec:real}

The topology of a complex slice is not, by itself, the correct count of
obstructions on a symmetric slice domain. Reflection imposes an additional
constraint.

Let $D\subset\C$ be connected, invariant under $\tau(z)=\bar z$, and meeting
$\R$. Suppose
\[
 g(x+Iy)=P(x,y)+I Q(x,y)
\]
is a smooth axial Fueter-regular function on its circularization, represented
on $D$ by smooth coefficients with
\begin{equation}
 P(x,-y)=P(x,y),\qquad Q(x,-y)=-Q(x,y).
 \label{eq:parity}
\end{equation}
The quotient $Q/y$ is interpreted by its smooth extension at $y=0$.
Define $\omega_g,\eta_g$ on all of $D$ by replacing $r$ with signed $y$ in
\eqref{eq:intro-omega}--\eqref{eq:intro-eta}.

\begin{theorem}[Reflection-invariant obstruction]\label{thm:reflection}
The two forms are smooth and closed on $D$, and
\begin{equation}
 \tau^*\omega_g=\omega_g,\qquad \tau^*\eta_g=\eta_g.
 \label{eq:invariance}
\end{equation}
There is a global slice-regular $f$ on the circularization of $D$ with
$Tf=g$ if and only if all their periods vanish. Thus the obstruction lies in
\begin{equation}
 H^1(D;\R)^{\tau^*}\otimes_\R(\HH\oplus\HH).
 \label{eq:invariant-cohomology}
\end{equation}
\end{theorem}
\begin{proof}
Closedness away from the real axis is the calculation in
Lemma~\ref{lem:closed}; smoothness extends it to the axis. The parity in
\eqref{eq:parity}, including the sign of $\dd y$ under reflection, gives
\eqref{eq:invariance}.

If all periods vanish, choose global potentials $C,D_0$ with
$\dd C=\omega_g$ and $\dd D_0=\eta_g$. The functions $C\circ\tau-C$ and
$D_0\circ\tau-D_0$ are constant. They vanish at any real basepoint, so both
potentials are even in $y$. Put
\[
 A=D_0+xC,\qquad B=yC.
\]
The proof of Lemma~\ref{lem:potentials}, with signed $y$, proves the
Cauchy--Riemann equations, including at $y=0$ by continuity. Also, $A$ is even
and $B$ odd, so $A+\ic B$ is a holomorphic stem. Its induced function is
smooth across the real axis and has Laplacian $g$. Necessity follows from the
same potential identities, with $B/y$ extended smoothly at $y=0$.
\end{proof}

\begin{theorem}[Which holes contribute]\label{thm:real-holes}
Assume $D$ has finite planar topology with a basis of counterclockwise hole
cycles, and its complementary components are permuted by conjugation. Let
$h$ be the number of pairs of distinct conjugate bounded complementary
components. Then the period obstruction map is onto $\HH^{2h}$ and
\begin{equation}
 \dim_{\HH}\bigl(\AM(D)/T\SR(D)\bigr)=2h.
 \label{eq:real-dimension}
\end{equation}
Every conjugation-fixed bounded complementary component contributes zero.
\end{theorem}
\begin{proof}
Reflection reverses planar orientation. If $\gamma_K$ is a counterclockwise
cycle around a hole $K$, then
$\tau_*[\gamma_K]=-[\gamma_{\bar K}]$. Invariance of the forms gives
\[
 \Per_g(\gamma_{\bar K})=-\Per_g(\gamma_K).
\]
For a fixed hole this forces its periods to be zero. A conjugate pair leaves
one freely specifiable pair of quaternions.

To realize the free classes, choose $p=u+\ic v$ in the upper member of each
pair and consider the intrinsic rational function
\begin{equation}
 k_p(z)=\frac{1}{2\pi\ic}
 \left(\frac1{z-p}-\frac1{z-\bar p}\right)
 =\frac{v}{\pi((z-u)^2+v^2)}.
 \label{eq:paired-derivative}
\end{equation}
Its local primitive can be written
\begin{equation}
 H_p(z)=\frac{1}{2\pi\ic}
 \log\frac{z-p}{z-\bar p},
 \label{eq:paired-log}
\end{equation}
up to a real constant and branch. Every local branch near a real point is
real on the real line, because the ratio in \eqref{eq:paired-log} has modulus
one there. Thus it induces a local slice-regular function across the axis.
The change around a loop is
$\Ind(\gamma,p)-\Ind(\gamma,\bar p)\in\mathbb Z$.
Consequently $T\cI(H_p)$ and $T\cI(zH_p)$ are globally defined smooth axial
functions on the circularization of $D$. They have, around the chosen upper
hole, periods $(0,1)$ and $(1,0)$, and opposite periods around its reflection.
Their right quaternionic multiples realize all classes. Exactness follows
from Theorem~\ref{thm:reflection}.
\end{proof}

\begin{remark}[Known arctangent kernels]
Locally, $H_p$ differs by a real constant from
$\pi^{-1}\arctan((z-u)/v)$. These are scaled and translated versions of the
arctangent primitives already used for spherical Cauchy kernels in
\cite{CSS,CPSS}. The new use here is their role as a normalized basis of the
global obstruction space, not their existence as special functions.
\end{remark}

\begin{example}[A real-centered annulus versus paired punctures]
For $D=\{z:a<|z|<b\}$, $0<a<b$, the bounded complementary component is
fixed by conjugation. Every smooth axial Fueter-regular function on the
quaternionic annulus $a<|q|<b$ therefore has a global slice-regular primitive.
The nonzero first homology of each complex slice does not obstruct inversion.
By contrast, deleting $p$ and $\bar p$ from a disk containing both and meeting
$\R$ produces a quaternionic domain with one deleted nonreal sphere and two
quaternionic obstruction coordinates.
\end{example}

\section{Stability and approximation with prescribed periods}
\label{sec:approx}

\subsection{Period estimates and a normalized inverse}
For $g=P+I Q$, define the profile norm
\[
 M_g(x,r)=\bigl(|P(x,r)|^2+|Q(x,r)|^2\bigr)^{1/2}.
\]
The sphere has area $4\pi$, and averaging over $I\in\SSS$ gives
\begin{equation}
 \frac1{4\pi}\int_{\SSS}|P+I Q|^2\dd\sigma(I)=|P|^2+|Q|^2.
 \label{eq:sphere-average}
\end{equation}
In particular, $M_g\le\sup_I|g(x+Ir)|$.

\begin{proposition}[Quantitative obstruction]\label{prop:estimates}
Let $\gamma\subset U$ have length $L_\gamma$ and
$R_\gamma=\max_{z\in\gamma}|z|$. Then
\begin{align}
 |a_\gamma|&\le\frac{L_\gamma}{2}\sup_\gamma M_g,\\
 |b_\gamma|&\le\frac{L_\gamma R_\gamma}{2}\sup_\gamma M_g.
 \label{eq:period-bounds}
\end{align}
If $h$ is a global Fueter image, and
$K_\gamma=\{x+Ir:x+\ic r\in\gamma,\ I\in\SSS\}$, then
\begin{equation}
 \sup_{K_\gamma}|g-h|
 \ge\frac{2\sqrt{|a_\gamma|^2+|b_\gamma|^2}}
 {L_\gamma\sqrt{1+R_\gamma^2}}.
 \label{eq:distance-bound}
\end{equation}
For a zero-period function, the inverse normalized by $C(z_0)=D(z_0)=0$
satisfies, along a path of length $L$ contained in $|z|\le R$,
\begin{equation}
 |f(x+Ir)|\le LR\sup_{\text{path}}M_g.
 \label{eq:inverse-bound}
\end{equation}
\end{proposition}
\begin{proof}
For a tangent vector $(\dot x,\dot r)$, Cauchy--Schwarz gives
\[
 |\omega_g(\dot\gamma)|\le\tfrac12 M_g|\dot\gamma|.
\]
The coefficients in $2\eta_g(\dot\gamma)$ are
$x\dot x+r\dot r$ and $r\dot x-x\dot r$; the sum of their squares is
$|z|^2|\dot\gamma|^2$. This proves the analogous bound with an additional
factor $|z|$ and hence \eqref{eq:period-bounds}.
Apply both estimates to $g-h$, whose periods equal those of $g$, and use
\eqref{eq:sphere-average} to get \eqref{eq:distance-bound}.
Finally $|C|\le LM/2$, $|D|\le LRM/2$, and $|q|\le R$ in
$f=qC+D$, proving \eqref{eq:inverse-bound}.
\end{proof}

\begin{corollary}[Closed range and continuous correction]\label{cor:closed}
The range $T\SR(U)$ is closed within $\AM(U)$ for locally uniform
convergence. For finite topology, $S\Pi_U$ and $1-S\Pi_U$ are continuous
projections onto the two summands in \eqref{eq:direct-sum}.
\end{corollary}
\begin{proof}
The range is the intersection of the kernels of all period functionals by
Theorem~\ref{thm:period}. Each functional is continuous by
Proposition~\ref{prop:estimates}. For finite topology, $S$ is a finite linear
combination of fixed smooth functions, hence continuous, and
$\Pi_U S=1$ gives the projection identities.
\end{proof}

\subsection{Rational approximation}
Under the hypotheses of Theorem~\ref{thm:split}, let
$E=\{p_1,\ldots,p_m,\bar p_1,\ldots,\bar p_m\}$.
Let $\mathcal R_E$ denote slice functions induced by rational $\HC$-valued
stems $R$ satisfying
\[
 R(\bar z)=\kappa(R(z)),
\]
with finite poles only in $E$. A pole at infinity is allowed. Equivalently,
a common denominator can be chosen to be a real polynomial with zeros in
$E$, and the numerator has right quaternionic coefficients. Thus these are
genuine rational stems, not independently chosen rational functions on each
quaternionic slice.

\begin{theorem}[Period-corrected Runge approximation]\label{thm:runge}
Every $g\in\AM(U)$ has a sequence $r_n\in\mathcal R_E$ such that
\begin{equation}
 T r_n+S\Pi_Ug\longrightarrow g
 \label{eq:runge}
\end{equation}
locally uniformly on $\Omega_U$, together with every fixed order of real
derivatives. The approximants in \eqref{eq:runge} have exactly the periods of
$g$. Furthermore,
\begin{equation}
 \overline{T\mathcal R_E}^{\,\mathrm{loc}}=T\SR(U)
 \quad\text{inside }\AM(U).
 \label{eq:runge-closure}
\end{equation}
\end{theorem}
\begin{proof}
By Theorem~\ref{thm:split}, write $g-S\Pi_Ug=Tf$ and let $F$ be the
holomorphic stem of $f$ on $U\cup\bar U$. The pole set $E\cup\{\infty\}$
meets each complementary component of this symmetric planar set. Scalar
Runge approximation \cite{Runge}, applied to the four complex components of
$F$, gives rational functions $R_n^0$ with these allowed poles converging to
$F$ on compact subsets. Replace them by
\[
 R_n(z)=\tfrac12\bigl(R_n^0(z)+\kappa(R_n^0(\bar z))\bigr).
\]
This is still rational and holomorphic off $E$, has the required symmetry,
and still converges locally uniformly to $F$. Choosing a real common
denominator shows that its numerator coefficients belong to $\HH$.

Cauchy estimates on slightly larger compact subsets give convergence of all
complex derivatives. The stem-to-slice operation and the axial formula for
$T$ then give convergence of all fixed real derivatives on compact subsets of
$\Omega_U$, where $r$ is bounded away from zero. Set $r_n=\cI(R_n)$.
All $Tr_n$ have zero periods, which proves the period assertion.
The inclusion from right to left in \eqref{eq:runge-closure} follows by
applying the construction to a zero-period $g$; the other follows from
Corollary~\ref{cor:closed}.
\end{proof}

The correction term is necessary, not merely a feature of this proof. A
nonzero period rules out approximation by \emph{any} sequence of global
Fueter images on a compact set containing its supporting cycle, by
\eqref{eq:distance-bound}. This remains true even if the approximating
slice-regular functions are not rational.

\section{An endpoint removable-potential lemma}
\label{sec:weak}

The singularity classification rests on a two-dimensional endpoint fact.
We include its proof because the difference between $L^1$ and weaker
integrability is decisive.

\begin{lemma}[A puncture does not obstruct an integrable gradient]
\label{lem:W11}
Let $B_R\subset\R^2$ be a disk centered at zero and
$u\in C^1(B_R\setminus\{0\};\R^N)$. If
$\nabla u\in L^1(B_R)$, then $u$ extends, up to its value at zero, to a
function in $W^{1,1}_{\mathrm{loc}}(B_R)$ whose weak gradient is its original
gradient off zero.
\end{lemma}
\begin{proof}
Choose $0<R_0<R$. In polar coordinates, radial integration gives
\[
 |u(\rho,\theta)|\le |u(R_0,\theta)|+
 \int_\rho^{R_0}|u_t(t,\theta)|\dd t.
\]
Multiplication by $\rho$ and integration in $0<\rho<R_0$ and $\theta$ show
that $u\in L^1(B_{R_0})$: the integral involving $u_t$ is at most
\[
 \frac12\int_0^{2\pi}\int_0^{R_0}t^2|u_t(t,\theta)|\dd t\dd\theta
 \le\frac{R_0}{2}\int_{B_{R_0}}|\nabla u|\dd x\dd y.
\]
Also, with $H(t)=\int_0^{2\pi}|u_t(t,\theta)|\dd\theta$,
\[
 \rho\int_0^{2\pi}|u(\rho,\theta)|\dd\theta
 \le O(\rho)+\int_\rho^{R_0}\frac\rho t\,tH(t)\dd t\longrightarrow0.
\]
The convergence follows by dominated convergence because $tH(t)$ is
integrable. Integrating by parts outside the disk of radius $\rho$, the
inner boundary term is bounded by this last expression times the supremum
of the test function, and therefore tends to zero. This proves the weak
derivative assertion on $B_{R_0}$. Since $R_0<R$ is arbitrary, the result
follows.
\end{proof}

\begin{lemma}[Weak Cauchy--Riemann removability]\label{lem:weakCR}
If $F\in W^{1,1}_{\mathrm{loc}}(B_R;\C^N)$ satisfies
$\partial_{\bar z}F=0$ in distributions, then $F$ agrees almost everywhere
with a holomorphic function on $B_R$.
\end{lemma}
\begin{proof}
Convolve with a smooth approximate identity on smaller disks. The
convolutions $F_\varepsilon$ are smooth, satisfy the Cauchy--Riemann equation,
and converge to $F$ in local $L^1$. For holomorphic functions, the mean-value
formula on disks gives, on a compact disk $K$ contained in a slightly larger
disk $K'$, an estimate
\[
 \sup_K|H|\le C_{K,K'}\norm{H}_{L^1(K')}.
\]
Apply this to differences of convolutions. They are uniformly Cauchy on
compact subsets and hence converge to a holomorphic limit, which agrees
almost everywhere with the $L^1$ limit $F$.
\end{proof}

\section{Complete classification of integrable spherical singularities}
\label{sec:singularities}

Fix $p=u+\ic v\in\C^+$ and $0<R<v$. Set
\begin{equation}
 T_R=\{q:\rho(q)<R\},\qquad
 \rho(q)=\sqrt{(\Rea q-u)^2+(|\Ima q|-v)^2}.
 \label{eq:tube}
\end{equation}
Thus $T_R$ is the circularization of $|z-p|<R$ and $S_p=\{\rho=0\}$.
The Euclidean volume element is
\begin{equation}
 \dd V=r^2\dd x\dd r\dd\sigma(I),\qquad
 \int_{\SSS}\dd\sigma=4\pi.
 \label{eq:volume}
\end{equation}

For an axial $g=P+I Q$, averaging $g$ and $-Ig$ gives $P$ and $Q$:
\[
 P=\frac1{4\pi}\int_{\SSS}g\dd\sigma,
 \qquad Q=-\frac1{4\pi}\int_{\SSS}I g\dd\sigma.
\]
Consequently, on a sufficiently small tube, four-dimensional $L^1$
integrability of $g$ is equivalent to planar $L^1$ integrability of $P,Q$.
The same equivalence holds for every $L^p$ with $1\le p<\infty$, by Jensen's
inequality and $|P+I Q|\le|P|+|Q|$; $r$ stays bounded above and away from zero.

For $g\in\AM(T_R\setminus S_p)$, define its two residues by any small
counterclockwise meridian circle:
\begin{equation}
 \operatorname{res}_{p,1}g=a=\int_{|z-p|=\rho}\omega_g,
 \qquad
 \operatorname{res}_{p,0}g=b=\int_{|z-p|=\rho}\eta_g.
 \label{eq:residues}
\end{equation}
Closedness makes them independent of $\rho\in(0,R)$.
The notation $\AM$ on a tube means the same axial class as before.

\begin{theorem}[$L^1$ spherical normal form]\label{thm:L1}
Suppose $g\in\AM(T_R\setminus S_p)$ is integrable on some smaller punctured
tube $T_{R_0}\setminus S_p$. Then
\begin{equation}
 g=G_{p,1}a+G_{p,0}b+h,
 \label{eq:normal-form}
\end{equation}
where $a,b$ are its residues and $h$ extends as a smooth axial
Fueter-regular function through all of $S_p$.
The coefficients are unique. In particular,
\begin{equation}
 g\text{ is removable across }S_p
 \quad\Longleftrightarrow\quad a=b=0.
 \label{eq:L1-removal}
\end{equation}
\end{theorem}
\begin{proof}
The explicit formulas \eqref{eq:G0}--\eqref{eq:G1} show that, near $p$,
\[
 |P_k|+|Q_k|=O\bigl(\rho^{-1}+|\log\rho|+1\bigr),\qquad k=0,1.
\]
They are therefore $L^1$ on a small tube. Subtract their residue-normalized
combination from $g$ and call the result $g_0=P_0^*+I Q_0^*$. By
Proposition~\ref{prop:normalization}, both periods of $g_0$ vanish.
Theorem~\ref{thm:period} supplies a single-valued holomorphic stem
$F=A+\ic B$ on the punctured disk such that $T\cI(F)=g_0$.

Let $C=B/r$ and $D=A-xB/r$. Lemma~\ref{lem:potentials} gives
\[
 \dd C=\omega_{g_0},\qquad \dd D=\eta_{g_0}.
\]
Since $P_0^*,Q_0^*$ are planar $L^1$ and $x,r$ are bounded on a smaller disk,
$\nabla C$ and $\nabla D$ are $L^1$. Apply Lemma~\ref{lem:W11} to all eight
real components. It follows that
\[
 A=D+xC,\qquad B=rC
\]
extend to $W^{1,1}_{\mathrm{loc}}$ functions across the puncture. Their weak
Cauchy--Riemann equations hold there because they hold almost everywhere off
the puncture, and their weak derivatives are the extended original ones.
Lemma~\ref{lem:weakCR} now extends $F$ holomorphically. Its induced slice
function and its Fueter image extend smoothly across the entire sphere.
This proves \eqref{eq:normal-form}.

A smooth extension has zero periods by Stokes' theorem on a meridian disk.
Thus the coefficients of any such decomposition must be exactly
\eqref{eq:residues}; uniqueness and \eqref{eq:L1-removal} follow.
\end{proof}

\begin{corollary}[Finite many integrable singular spheres]\label{cor:many}
Let an axial Fueter-regular function be defined on an axial domain with
finitely many pairwise disjoint nonreal spheres removed, and be locally
$L^1$ near each removed sphere. Subtracting one residue-normalized pair of
local generators at each sphere leaves a function that extends across every
removed sphere, wherever those generators are defined on the ambient domain.
On domains meeting the real axis one may use the paired generators
$T\cI(zH_p),T\cI(H_p)$ instead.
\end{corollary}
\begin{proof}
Generators based at other spheres are smooth near a given sphere. The
paired and unpaired generators differ, near the upper puncture, by a smooth
Fueter image because the reflected pole lies outside that disk. Apply
Theorem~\ref{thm:L1} at each sphere. The extensions agree with the original
function off the spheres and hence fit together.
\end{proof}

\begin{proposition}[Sharpness of the $L^1$ hypothesis]\label{prop:L1sharp}
For every $0<s<1$, there is an axial Fueter-regular function locally in
$L^s$ near $S_p$ that has $a=b=0$ but is not removable.
\end{proposition}
\begin{proof}
Take the single-valued holomorphic stem $F(z)=(z-p)^{-1}$ on the punctured
disk and put $g=T\cI(F)$. Its periods vanish because it is a global Fueter
image there. The derivative $F'=-1/(z-p)^2$, together with
\eqref{eq:fueter-formula}, gives $|g|=O(\rho^{-2})$ uniformly in $I$.
Thus $g\in L^s$ for $s<1$, since
$\int_0^{R_0}\rho^{1-2s}\dd\rho<\infty$.

It is not removable. Otherwise the local inverse theorem on the full
meridian disk would give a holomorphic stem $F_0$ with
$T\cI(F_0)=g$. The kernel calculation in Theorem~\ref{thm:period} would make
$F-F_0=za+b$ on the punctured disk for fixed $a,b\in\HH$, impossible because
$F$ has a pole. This argument also avoids relying on a pointwise lower bound
in a possibly exceptional quaternionic direction.
\end{proof}

\section{Residues, exact energy, and the distributional source}
\label{sec:energy}

\subsection{The uniform leading term}
\begin{lemma}[Asymptotic profile]\label{lem:leading}
Under the hypotheses of Theorem~\ref{thm:L1}, write
$z=p+\rho e^{\ic\theta}$ and set
\begin{equation}
 c_0=ua+b,\qquad c_1=va.
 \label{eq:cs}
\end{equation}
Uniformly in $\theta$ and $I\in\SSS$,
\begin{equation}
 g(x+Ir)=\frac{e^{-I\theta}(c_0+I c_1)}{\pi v\rho}
             +O(1+|\log\rho|).
 \label{eq:leading}
\end{equation}
The quaternionic coefficients remain on the right; no commutation with $I$
is implicit.
\end{lemma}
\begin{proof}
For a local stem primitive of the singular part, use
\[
 F(z)=(za+b)\ell_p(z).
\]
Its derivative is
\[
 F'(z)=a\ell_p(z)+\frac{za+b}{2\pi\ic(z-p)}
      =\frac{c_0+\ic c_1}{2\pi\ic\rho e^{\ic\theta}}
          +O(1+|\log\rho|).
\]
A branch with $\theta$ in any interval of length $2\pi$ gives the stated
uniform error; changing the branch adds an affine kernel and does not change
$g$. If $F'=U+\ic V$ and $F=A+\ic B$, then
\[
 P=-\frac{2V}{r},\qquad Q=\frac{2U}{r}-\frac{2B}{r^2}.
\]
Here $B=O(1+|\log\rho|)$, and $r=v+O(\rho)$.
Since $(\ic e^{\ic\theta})^{-1}=-\sin\theta-\ic\cos\theta$, the leading
coefficients are
\begin{align}
 P_*&=\frac{\cos\theta\,c_0+\sin\theta\,c_1}{\pi v\rho},\\
 Q_*&=\frac{-\sin\theta\,c_0+\cos\theta\,c_1}{\pi v\rho}.
 \label{eq:leading-PQ}
\end{align}
Their combination $P_*+I Q_*$ is exactly the first term of
\eqref{eq:leading}. The removable part is bounded on a smaller tube.
\end{proof}

\subsection{Energy and the sharp \texorpdfstring{$L^2$}{L2} threshold}
\begin{theorem}[Exact logarithmic energy and finite renormalized limit]
\label{thm:energy}
Let $g$ satisfy Theorem~\ref{thm:L1}, and fix a sufficiently small $R>0$.
Then the limit
\begin{equation}
 \lim_{\varepsilon\downarrow0}\left[
 \int_{T_R\setminus\overline{T_\varepsilon}}|g|^2\dd V
 -8\bigl(|ua+b|^2+v^2|a|^2\bigr)\log\frac R\varepsilon
 \right]
 \label{eq:renormalized}
\end{equation}
exists and is finite. In particular the coefficient of logarithmic divergence
is positive exactly when the singularity is nonremovable.
\end{theorem}
\begin{proof}
Let $g_*=P_*+I Q_*$ be the first term of \eqref{eq:leading} and
$S=|c_0|^2+|c_1|^2$. The real orthogonal rotation in
\eqref{eq:leading-PQ} gives
\[
 |P_*|^2+|Q_*|^2=\frac{S}{\pi^2v^2\rho^2}.
\]
By \eqref{eq:sphere-average} and \eqref{eq:volume}, with
$r=v+\rho\sin\theta$,
\begin{align*}
 \int_{T_R\setminus\overline{T_\varepsilon}}|g_*|^2\dd V
 &=\frac{4S}{\pi v^2}\int_\varepsilon^R\frac1\rho
     \int_0^{2\pi}(v+\rho\sin\theta)^2\dd\theta\dd\rho\\
 &=8S\log\frac R\varepsilon+
     \frac{2S}{v^2}(R^2-\varepsilon^2).
\end{align*}
For the remainder $E=g-g_*$, Lemma~\ref{lem:leading} gives
$|E|\le C(1+|\log\rho|)$. Hence
\[
 \bigl||g|^2-|g_*|^2\bigr|
 \le2|g_*||E|+|E|^2
\]
is integrable down to $\rho=0$: after including the bounded $r^2$ and polar
Jacobian $\rho$, the relevant radial bounds are
$1+|\log\rho|$ and $\rho(1+|\log\rho|)^2$.
The difference of the two energy integrals therefore has a finite limit.
Since $S=|ua+b|^2+v^2|a|^2$, this proves \eqref{eq:renormalized}.
Finally $v>0$ implies $S=0$ exactly when $a=b=0$.
\end{proof}

\begin{corollary}[Sharp $L^p$ classification]\label{cor:Lp}
For locally axial Fueter-regular germs along $S_p$:
\begin{enumerate}[label=\textup{(\roman*)}]
\item if $1\le s<2$, the locally $L^s$ germs modulo removable germs form a
right quaternionic vector space of dimension two, with coordinates $(a,b)$;
\item if $2\le s\le\infty$, every locally $L^s$ germ is removable;
\item for every $0<s<1$, the residues are not a complete invariant modulo
removable germs.
\end{enumerate}
\end{corollary}
\begin{proof}
The bounds on the two generators show they belong to $L^s$ for $s<2$.
For $s\ge1$, local $L^s$ implies local $L^1$ on a finite-volume tube, so the
normal form applies. For $s\ge2$, it also implies local $L^2$, and
Theorem~\ref{thm:energy} forces both residues to vanish. The last assertion
is Proposition~\ref{prop:L1sharp}.
\end{proof}

These are two different sharp thresholds. At $L^1$ the periods become a
\emph{complete} singularity invariant; at $L^2$ they must vanish. The first
threshold is not a consequence of the second.

\subsection{The affine surface source}
For a smooth surface $S$, the notation $j\delta_S$ means the
quaternion-valued distribution
\[
 \langle j\delta_S,\varphi\rangle=\int_S j(q)\varphi(q)\dd S(q)
\]
for real scalar test functions $\varphi$. The first-order operator $\D$ acts
with quaternionic units on the left, also in the distributional definition.

\Needspace{17\baselineskip}
\begin{theorem}[Distributional defect and layer-potential normal form]
\label{thm:defect}
Extend an $L^1$ function $g$ from Theorem~\ref{thm:L1} arbitrarily on the
measure-zero sphere $S_p$. Then, on a sufficiently small full tube,
\begin{equation}
 \D g=j_{a,b}\,\delta_{S_p},\qquad
 j_{a,b}(q)=\frac2v(qa+b).
 \label{eq:defect}
\end{equation}
Let
\[
 E(q)=\frac{\bar q}{2\pi^2|q|^4},\qquad q\ne0.
\]
Then equivalently, modulo a smooth axial Fueter-regular function,
\begin{equation}
 g(q)=\frac2v\int_{S_p}E(q-s)(sa+b)\dd S(s).
 \label{eq:layer}
\end{equation}
The density is affine in the source point $s$, with its coefficients on the
right of the Cauchy--Fueter kernel.
\end{theorem}
\begin{proof}
The outward unit normal of the small tube $T_\rho$ at
$x=u+\rho\cos\theta$, $r=v+\rho\sin\theta$ is
\[
 n_\rho=\cos\theta+I\sin\theta=e^{I\theta},
 \qquad \dd S_{\partial T_\rho}=r^2\rho\dd\theta\dd\sigma(I).
\]
Integration by parts on the complement of $T_\rho$, with a compactly
supported real test function, gives
\[
 \langle\D g,\varphi\rangle
 =\lim_{\rho\downarrow0}\int_{\partial T_\rho}n_\rho g\varphi\dd S.
\]
The sign is positive: the inner boundary normal of the punctured region is
$-n_\rho$. Multiplying \eqref{eq:leading} on the left by $n_\rho$ cancels
$e^{-I\theta}$. The error contributes
$O(\rho(1+|\log\rho|))$, which tends to zero. The limit is
\begin{align*}
 \frac{v}{\pi}\int_{\SSS}\int_0^{2\pi}
 (c_0+I c_1)\varphi(u+Iv)\dd\theta\dd\sigma(I)
 &=2v\int_{\SSS}(c_0+I c_1)\varphi(u+Iv)\dd\sigma(I)\\
 &=\frac2v\int_{S_p}(qa+b)\varphi(q)\dd S(q).
\end{align*}
This proves \eqref{eq:defect}.

For completeness, $\D E=0$ away from zero, by direct differentiation, and
$\D E=\delta_0$ distributionally. Indeed, on a radius-$\varepsilon$
three-sphere, $nE=1/(2\pi^2\varepsilon^3)$ and its integral is one; the
same integration-by-parts argument proves the distributional identity.
The layer potential in \eqref{eq:layer} is locally integrable and smooth off
$S_p$. Local integrability follows from the integrability of $|q-s|^{-3}$
in four dimensions and Fubini over the compact source sphere. Its Dirac
image is exactly $j_{a,b}\delta_{S_p}$.

The layer is axial. For densities $1$ and $s$, conjugation by a unit
quaternion rotates the integration sphere and transforms the resulting
function by the same conjugation. The stabilizer of $I$ therefore forces
its value at $x+Ir$ to lie in $\R+I\R$; the two real coefficients depend
only on $x,r$. Right multiplication by $a,b$ preserves the axial form.

Subtract the layer from $g$. The difference is locally integrable and is
annihilated by $\D$ in distributions across the sphere. Applying
$\overline\D$ makes every component weakly harmonic, so the difference is
smooth by the harmonic Weyl lemma. One may prove this lemma by mollification:
mollified harmonic functions are controlled on smaller balls by their local
$L^1$ norms, and hence converge locally uniformly to a harmonic
representative. The distributional Dirac equation then holds classically.
The extension remains axial by continuity. This proves \eqref{eq:layer}.
\end{proof}

\begin{corollary}[Geometric form of the energy coefficient]
\label{cor:source-energy}
With $j_{a,b}$ as in \eqref{eq:defect},
\begin{equation}
 8\bigl(|ua+b|^2+v^2|a|^2\bigr)
 =\frac1{2\pi}\int_{S_p}|j_{a,b}(q)|^2\dd S(q).
 \label{eq:source-energy}
\end{equation}
Thus the logarithmic energy coefficient is the squared $L^2$ norm of the
surface source divided by $2\pi$.
\end{corollary}
\begin{proof}
The sphere average of $|c_0+I c_1|^2$ is $|c_0|^2+|c_1|^2$, and
$\dd S=v^2\dd\sigma$. Since $j_{a,b}=2(c_0+I c_1)/v$, its squared norm
integrates to $16\pi(|c_0|^2+|c_1|^2)$.
\end{proof}

\begin{corollary}[The ordinary total flux is incomplete]
\label{cor:flux}
For every sufficiently small tube enclosing $S_p$,
\begin{equation}
 \int_{\partial T_\rho}n_\rho g\dd S=8\pi v(ua+b).
 \label{eq:flux}
\end{equation}
In particular a nonremovable $L^1$ singularity can have zero total
Cauchy--Fueter flux: choose $a\ne0$ and $b=-ua$.
\end{corollary}
\begin{proof}
The boundary integral is independent of $\rho$ by the divergence theorem on
an annular tube. Its small-radius limit is the total mass of
\eqref{eq:defect}. The average of $I$ over $\SSS$ vanishes, yielding
\eqref{eq:flux}. If $a\ne0$ and $b=-ua$, the second coefficient $c_1=va$
is nonzero, so Theorem~\ref{thm:L1} and Theorem~\ref{thm:energy} show that the
singularity is nevertheless nonremovable.
\end{proof}

\begin{corollary}[Two conserved fluxes recover both residues]
\label{cor:twoflux}
Put $H_u(q)=3(\Rea q-u)+\Ima q$ and define
\begin{equation}
 \Phi_0(\rho)=\int_{\partial T_\rho}n_\rho g\dd S,
 \qquad
 \Phi_1(\rho)=\int_{\partial T_\rho}H_u(q)n_\rho g\dd S.
 \label{eq:twoflux}
\end{equation}
Both quantities are independent of $\rho$ on sufficiently small tubes, and
\begin{equation}
 a=-\frac{\Phi_1}{8\pi v^3},\qquad
 b=\frac{\Phi_0}{8\pi v}-ua.
 \label{eq:recover-flux}
\end{equation}
Thus removability in the $L^1$ axial class is equivalent to simultaneous
vanishing of these two conserved three-dimensional boundary integrals.
The energy coefficient can also be written as
\begin{equation}
 \frac1{8\pi^2}
 \left(\frac{|\Phi_0|^2}{v^2}+\frac{|\Phi_1|^2}{v^4}\right).
 \label{eq:flux-energy}
\end{equation}
\end{corollary}
\begin{proof}
The weight $H_u$ is right Fueter regular:
$\sum_{\mu=0}^3(\partial_{x_\mu}H_u)e_\mu=3-3=0$,
where $e_0=1$, $e_1=\ii$, $e_2=\jj$, and $e_3=\kk$.
For a smooth right-regular weight $H$, the componentwise divergence theorem
therefore gives, on a region where $\D g=0$,
\[
 \int_{\partial V}Hn g\dd S
 =\int_V\left(\sum_{\mu=0}^3(\partial_{x_\mu}H)e_\mu\right)g\dd V
   +\int_VH\D g\dd V=0.
\]
Apply this to annular tubes, first with $H=1$ and then with $H=H_u$.
The ordinary flux is $\Phi_0=8\pi v c_0$ by
Corollary~\ref{cor:flux}. To compute the weighted flux, take the small-tube
limit in the proof of Theorem~\ref{thm:defect}. On the source sphere,
$H_u(u+Iv)=Iv$, and therefore
\[
 \Phi_1=2v\int_{\SSS}Iv(c_0+I c_1)\dd\sigma(I)
       =-8\pi v^2c_1=-8\pi v^3a.
\]
This proves \eqref{eq:recover-flux}. Substitution into
$8(|c_0|^2+v^2|a|^2)$ gives \eqref{eq:flux-energy}.
\end{proof}

\begin{remark}[Covariance under a real translation]
Under $q'=q-t$ with real $t$, the slope residue remains $a$ and the intercept
becomes $b'=b+ta$. The sphere center becomes $u'=u-t$. Thus
$u'a+b'=ua+b$, and the energy coefficient is unchanged. The two coordinates
are affine-monodromy coordinates, not two independently translation-invariant
scalars.
\end{remark}

\section{Examples and reproducible checks}
\label{sec:checks}

\subsection{A singularity invisible to total flux}
Let $p=\ic$, so $S_p=\SSS$ is the unit imaginary sphere, and take
$g=G_{p,1}$. Its residues are $(a,b)=(1,0)$. It is locally $L^s$ for every
$s<2$, but it is not removable and has no single-valued slice-regular
primitive on a surrounding punctured tube. Nevertheless,
\[
 \int_{\partial T_\rho}n_\rho g\dd S=0.
\]
Its source density is $j(q)=2q=2I$ on the sphere, and
\[
 \int_{\varepsilon<\rho(q)<R}|g(q)|^2\dd V
 =8\log\frac R\varepsilon+O(1).
\]
The source has zero total mass but nonzero $L^2$ norm. For a version smooth
on the entire real axis, replace the local logarithmic generator with
$T\cI(zH_p)$ from \eqref{eq:paired-log}; the residues and the local energy
coefficient are the same.

\subsection{A nonremovable global Fueter image below \texorpdfstring{$L^1$}{L1}}
The example $g=T\cI((z-p)^{-1})$ has a globally single-valued primitive on
the punctured meridian and zero periods, but still fails to extend through
the sphere. This separates two issues that coincide only after imposing the
$L^1$ hypothesis: the existence of a global primitive, and removability of
the target singularity.

\subsection{What the verification script checks}
The archive includes \texttt{verify.py}, its JSON output, and a plain-text
summary. The script uses SymPy for exact symbolic identities and NumPy for
deterministic periodic quadrature. It performs thirteen symbolic checks:
both Vekua equations for both generators, closedness of both forms for both
generators, the leading angular energy identity, and all four potential
derivative identities for a holomorphic polynomial with genuinely complex
stem coefficients.

For $p=0.4+1.7\ic$, it integrates the two periods of both generators on
circles of radii $0.15$, $0.4$, and $0.9$, using $16\,384$ midpoint nodes.
The largest absolute discrepancy from $(0,1)$ and $(1,0)$ in the recorded
run was $1.11\times10^{-16}$. This is a floating-point consistency check,
not a rigorous error bound for numerical quadrature.
The script also checks the two conserved boundary fluxes in
Corollary~\ref{cor:twoflux} at the same radii, with non-real quaternionic
coefficients. The spherical average is evaluated analytically, and the
meridian integral uses $32\,768$ nodes; the largest component error in the
recorded run was $2.01\times10^{-12}$.

The energy check uses the non-real quaternionic coefficients
\[
 a=1+0.2\ii-0.3\jj+0.7\kk,
 \qquad b=-0.4+0.5\ii+0.9\jj-1.1\kk.
\]
For the same $p$, the predicted coefficient is
\[
 8\bigl(|0.4a+b|^2+1.7^2|a|^2\bigr)=50.392.
\]
The numerical quantity is the logarithmic shell energy
\begin{equation}
 J(\rho)=4\pi\rho^2\int_0^{2\pi}r^2
      \bigl(|P(x,r)|^2+|Q(x,r)|^2\bigr)\dd\theta,
 \label{eq:numerical-energy}
\end{equation}
where $x=u+\rho\cos\theta$, $r=v+\rho\sin\theta$ and
$g=G_{p,1}a+G_{p,0}b$. It converges to the coefficient in
Theorem~\ref{thm:energy}.
\begin{center}
\begin{tabular}{@{}rrr@{}}
\toprule
$\rho$ & $J(\rho)$ & Absolute discrepancy from $50.392$\\
\midrule
$10^{-1}$ & $50.44965278839$ & $5.76528\times10^{-2}$\\
$10^{-2}$ & $50.41030937149$ & $1.83094\times10^{-2}$\\
$10^{-3}$ & $50.39254549030$ & $5.45490\times10^{-4}$\\
$10^{-4}$ & $50.39201092790$ & $1.09279\times10^{-5}$\\
$3\times10^{-5}$ & $50.39200130732$ & $1.30732\times10^{-6}$\\
\bottomrule
\end{tabular}
\end{center}
No formal proof assistant was used. Symbolic computations verify finite
algebraic identities; the theorems rest on the analytical and topological
proofs in the preceding sections.

\section{Literature comparison and remaining scope}
\label{sec:literature}

\subsection{Comparison with the closest sources}
The foundational stem construction is due to Ghiloni--Perotti
\cite{GP}; it is used here only as a language for coordinating slices.
Colombo--Sabadini--Sommen \cite{CSS} developed the inverse Fueter mapping
in integral form, including the quaternionic axial case. Their inverse
construction, and its successors, establish the background local
surjectivity problem; we do not claim that problem as new.

Colombo--Pe\~na Pe\~na--Sabadini--Sommen \cite{CPSS} give explicit inverse
formulas on rectangular meridians, identify the polynomial kernel, and
compute spherical Cauchy kernels using arctangent primitives. In their
quaternionic degree-zero case, the kernel is the affine family used here.
The formulas $H_p$ and $zH_p$ are therefore closely related to known kernels.
The finite-hole exact sequence and the $L^1$ residue--defect--energy
classification are not statements we located in that paper.

Dong--Qian \cite{DQ} discuss axiality, surjectivity, and monomial versions of
the Fueter map, with integral formulas and kernel primitives. Broad
surjectivity terminology should not be interpreted, without examining domains
and branches, as a guarantee of a single-valued primitive on every
multiply connected meridian. Our examples demonstrate directly why a
period condition is indispensable for the latter assertion. We are not
asserting that the present theorem contradicts the properly localized or
branch-dependent content of those constructions.

De Martino--Diki--Guzm\'an Ad\'an \cite{DDG} relate Fueter maps to generalized
Cauchy--Kovalevskaya extension. De Martino--Pinton \cite{DP2026} study
factorizations and inversion problems for harmonic and second-order
polyanalytic axial classes. These are related but different function spaces
and inversion questions. Their work is relevant context rather than a premise
of the proofs above.

The ordinary Runge theorem \cite{Runge} is the only approximation machinery
used. Our approximation result adds the finite-dimensional period correction
and proves its necessity. The proof uses no new scalar rational
approximation theorem.

\subsection{Limits of the novelty check}
The search used combinations of ``inverse Fueter'', ``periods'',
``monodromy'', ``multiply connected'', ``cohomology'', and spherical
removability, followed by inspection of the closest primary sources.
Several queries produced many irrelevant results; absence from those results
is not evidence of absence from the literature. The assessment of novelty is
therefore deliberately limited to the exact formulations in the sources
actually inspected. A full bibliographic review could reveal an equivalent
cohomological statement, a differently normalized residue theorem, or a general
elliptic result containing a special case of the energy law.

This limitation concerns priority, not the logical status of the arguments
presented here. The paper supplies proofs of its assertions, explicitly
acknowledges the old local inverse and old kernels, and does not relabel a
standard spherical kernel as a discovery.

\subsection{What remains outside the proved results}
Three extensions are natural but are not needed for the conclusions.
For infinitely connected meridians, the zero-period criterion remains valid;
a useful explicit continuous splitting with infinitely many charges would
require convergence conditions and a choice of topology. In higher odd
Clifford dimensions, the Fueter kernel is a larger polynomial space, suggesting
polynomial rather than affine monodromy; extracting a comparably elementary
family of direct closed forms is a separate task. Finally, nonaxial
Fueter-regular functions can have much richer surface-supported source
families, so the two-residue classification here must not be transferred to
them without additional hypotheses.

\section{Conclusion}
The passage from local to global Fueter inversion is governed by two direct
period forms. On a meridian with $m$ holes they give an explicit cokernel
$\HH^{2m}$, a constructive correction operator, and the exact obstruction to
approximation by global Fueter images. Reflection reduces this count on
real-axis domains: fixed holes contribute no period obstruction, while
conjugate pairs do.

At an isolated nonreal sphere, the same two periods become a complete
singularity invariant precisely under the $L^1$ hypothesis proved here.
Their affine monodromy $qa+b$ is also the density, up to $2/v$, of the
sphere-supported Dirac defect. Its squared norm fixes the coefficient of
logarithmic energy divergence. The resulting distinction between a zero
\emph{total flux} and a zero \emph{source density} explains why one
quaternionic flux residue is not enough for a spherical singularity.

\clearpage
\phantomsection
\addcontentsline{toc}{section}{References}
\begin{thebibliography}{10}
\bibitem{GP}
R.~Ghiloni and A.~Perotti,
\emph{Slice regular functions on real alternative algebras},
Advances in Mathematics \textbf{226} (2011), 1662--1691.
\href{https://arxiv.org/abs/1008.4318}{arXiv:1008.4318}.

\bibitem{CSS}
F.~Colombo, I.~Sabadini, and F.~Sommen,
\emph{The inverse Fueter mapping theorem},
Communications on Pure and Applied Analysis \textbf{10} (2011), 1165--1181.
\href{https://doi.org/10.3934/cpaa.2011.10.1165}{doi:10.3934/cpaa.2011.10.1165}.

\bibitem{CPSS}
F.~Colombo, D.~Pe\~na Pe\~na, I.~Sabadini, and F.~Sommen,
\emph{A new integral formula for the inverse Fueter mapping theorem},
Journal of Mathematical Analysis and Applications \textbf{417} (2014), 112--122.
\href{https://arxiv.org/abs/1302.0685}{arXiv:1302.0685}.
In particular, Theorem~3, Corollary~1, and Example~2 of the preprint.

\bibitem{DQ}
B.~Dong and T.~Qian,
\emph{Further aspects of the Fueter mapping theorem},
preprint (2018).
\href{https://arxiv.org/abs/1805.02966}{arXiv:1805.02966}.

\bibitem{DDG}
A.~De Martino, K.~Diki, and A.~Guzm\'an Ad\'an,
\emph{On the connection between the Fueter--Sce--Qian theorem and the
 generalized CK-extension}, preprint (2022).
\href{https://arxiv.org/abs/2203.03490}{arXiv:2203.03490}.

\bibitem{DP2026}
A.~De Martino and S.~Pinton,
\emph{A variation of the inverse Fueter theorem and the generalized
 polyanalytic Cauchy--Kovalevskaya extension of order 2},
preprint, version 1, 7 August 2026.
\href{https://arxiv.org/abs/2608.07711}{arXiv:2608.07711}.

\bibitem{Runge}
C.~Runge,
\emph{Zur Theorie der eindeutigen analytischen Functionen},
Acta Mathematica \textbf{6} (1885), 229--244.
\href{https://doi.org/10.1007/BF02400416}{doi:10.1007/BF02400416}.

\bibitem{Supplied}
\emph{Classical Complex Theorems over the Quaternions:
 Slice-regular and Cauchy--Fueter analogues, with proofs},
supplied expository manuscript, 20 September 2026,
especially the section
``Fueter's mapping theorem and a constructive local inverse.''
Unpublished background supplied with the present task.
\end{thebibliography}
\end{document}
